% Chapitre « Inférence et disponibilité » (partie 5, sous-étape 5.11, jalons J7 et J8).
% Sources : ADR-0009, ADR-0010, ADR-0011, ADR-0040, ADR-0041, ADR-0042 ; code de
% src/foot_predictor/inference/ (rows.py, models.py, availability.py, predict.py, live_sources.py)
% et de src/foot_predictor/modeling/ (distributions.py, metrics.py). En cas d'écart, le code fait foi.

\chapter{Inférence et disponibilité}
\label{chap:inference}

Les chapitres précédents évaluent des modèles sur des saisons entières. Ce chapitre décrit ce qui se
passe pour \emph{un} match : calculer ses variables à la date de référence, décider s'il peut être
prédit, produire la loi du total et l'intervalle annoncé. Deux exigences guident tout : la
\emph{même fonction} qu'à l'entraînement (aucune formule propre à l'inférence) et \emph{aucune valeur
inventée} (une variable absente bloque la prédiction).

% ---------------------------------------------------------------------------------------------
\section{Variables d'un match quelconque}
\label{sec:inference-variables}

Notons \(\Phi\) la fonction qui construit le jeu d'entraînement (\code{features.dataset.build\_frame}) :
elle reçoit une table de matchs et renvoie une ligne par couple (match, équipe). Pour prédire les
matchs \(C_J\) d'un jour \(J\), on l'applique à la table
\[
  \hist_J = \{ m : d(m) \le J \}
  \quad \text{où les matchs de } C_J \text{ sont marqués joués, score } 0\text{-}0, \text{ sans statistiques,}
\]
puis on garde les lignes de \(C_J\) et l'on retire les colonnes d'après-match (buts, source des tirs).

\paragraph{Pourquoi c'est exact.} L'Elo d'avant-match du jour \(J\) n'utilise que les matchs terminés
avant \(J\) (les variations d'un jour s'appliquent ensemble, en fin de jour) ; les moyennes glissantes
et le calendrier ne lisent que les jours strictement antérieurs. Pour toute variable \(x_v\) d'un
match de \(C_J\), \(x_v\big(\hist_J^{(s)}\big)\) ne dépend donc pas du score fictif \(s\) attribué aux
matchs de \(C_J\), ni d'un match postérieur à \(J\). Les tests le vérifient en changeant ces scores ou
ces matchs, et en comparant les lignes obtenues à celles du jeu d'entraînement : identiques au bit
près sur 17 jours et 622 lignes.

\paragraph{Un jour à la fois.} Calculer deux jours \(J < J'\) d'un coup ferait entrer le score fictif
d'un match de \(J\) dans les variables de \(J'\) (l'Elo de fin de \(J\), la moyenne glissante de
\(J'\)). On applique donc \(\Phi\) jour par jour, avec un cache par (jour, version des données).

\paragraph{Modèle de rejeu.} Pour une date de la saison \(S\), le modèle est celui du pli \(S\) :
\(\hat\theta_S\) appris sur les saisons 2015-16 à \(S - 1\) par la procédure de l'évaluation, jamais
un modèle qui a vu la saison \(S\) (ADR-0011, règle 1). Ses prédictions reproduisent celles de
l'évaluation à la précision machine (écart maximal de \(\lambda\) de l'ordre de \(10^{-15}\)).

% ---------------------------------------------------------------------------------------------
\section{Loi du total et quantités publiées}
\label{sec:inference-loi}

Le modèle M3 donne, pour chaque équipe, \(\lambda_e = \exp(x_e^\top \hat\beta)\). Sous l'indépendance
conditionnelle des deux équipes (chapitre~\ref{chap:modeles}), le total suit
\(T \sim \mathcal{P}(\lambda_{\mathrm{dom}} + \lambda_{\mathrm{ext}})\), replié sur
\(\{0, \dots, 9, \ge 10\}\) :
\[
  \hat p_k = e^{-\lambda}\frac{\lambda^k}{k!} \ (k \le 9),
  \qquad
  \hat p_{10} = 1 - \sum_{k=0}^{9} \hat p_k,
  \qquad \lambda = \lambda_{\mathrm{dom}} + \lambda_{\mathrm{ext}} .
\]
On publie \(E[T] = \lambda\) (espérance de la loi non repliée), \(\hat P(T > 2{,}5) = 1 - \hat p_0 -
\hat p_1 - \hat p_2\), la loi \((\hat p_k)\) et l'intervalle de la section suivante.

% ---------------------------------------------------------------------------------------------
\section{Intervalle et couverture annoncée}
\label{sec:inference-intervalle}

Avec \(\hat F(k) = \sum_{j \le k} \hat p_j\), les bornes sont les plus petits entiers qui atteignent
les niveaux 10~\% et 90~\% :
\[
  q_{10} = \min\{k : \hat F(k) \ge 0{,}10\},
  \qquad
  q_{90} = \min\{k : \hat F(k) \ge 0{,}90\},
\]
la borne 10 signifiant \og 10 et plus \fg{}. La loi étant discrète, \([q_{10} ; q_{90}]\) ne couvre pas
exactement 80~\% ; on publie sa \emph{couverture annoncée}
\[
  \hat c = \hat P\big(T \in [q_{10} ; q_{90}]\big) = \hat F(q_{90}) - \hat F(q_{10} - 1) .
\]
Comme \(\hat F(q_{90}) \ge 0{,}90\) et \(\hat F(q_{10} - 1) < 0{,}10\), on a toujours \(\hat c > 0{,}80\) :
annoncer \og 80~\% \fg{} serait faux, d'où l'obligation de publier \(\hat c\) (ADR-0009). Le contrôle
compare la couverture observée à la moyenne des \(\hat c\), à \(\pm 3\) points.

\paragraph{Exemple.} Pour \(\lambda_{\mathrm{dom}} = 1{,}6\) et \(\lambda_{\mathrm{ext}} = 1{,}2\) :

\begin{center}
\begin{tabular}{lccccccccccc}
\toprule
\(k\) & 0 & 1 & 2 & 3 & 4 & 5 & 6 & 7 & 8 & 9 & \(\ge 10\) \\
\midrule
\(\hat p_k\) & 0,061 & 0,170 & 0,238 & 0,222 & 0,156 & 0,087 & 0,041 & 0,016 & 0,006 & 0,002 & 0,001 \\
\(\hat F(k)\) & 0,061 & 0,231 & 0,470 & 0,692 & 0,848 & 0,935 & 0,976 & 0,992 & 0,998 & 0,999 & 1 \\
\bottomrule
\end{tabular}
\end{center}

\(E[T] = 2{,}8\), \(\hat P(T > 2{,}5) = 0{,}531\), \(q_{10} = 1\), \(q_{90} = 5\) et
\(\hat c = 0{,}935 - 0{,}061 = 0{,}874\) : l'intervalle \([1 ; 5]\) est annoncé à 87,4~\%.

% ---------------------------------------------------------------------------------------------
\section{Matrice de disponibilité}
\label{sec:inference-disponibilite}

Pour chaque variable requise \(v\) du modèle et chaque équipe \(e\) d'un match du jour \(J\), un statut :

\begin{itemize}
  \item \textbf{présente} : la valeur \(x_{v,e}\) est calculée ;
  \item \textbf{manquante} : \(\Phi\) la laisse vide, avec la raison (aucun match de championnat dans les
        730 jours, aucun tir connu, équipe inconnue du référentiel) ;
  \item \textbf{périmée} : variable d'historique dont la source n'est complète que jusqu'au jour
        \(C < J - 1\) ; il manque peut-être des matchs antérieurs à \(J\) (ADR-0011, règle 5).
\end{itemize}

Le match est \code{available} si toutes ses variables sont présentes, \code{unavailable} sinon (raisons
listées, aucune prédiction) ; \code{out\_of\_scope} hors du top 5 ou hors saison régulière ;
\code{excluded} pour un match annulé ou abandonné ; \code{h2\_unavailable} pour l'horizon avec
composition. \emph{Aucune valeur de remplacement} : remplacer une variable par une moyenne ou par
zéro produirait un \(\lambda\) faux présenté avec la même confiance qu'un vrai.

% ---------------------------------------------------------------------------------------------
\section{Le live : sources superposées et fraîcheur}
\label{sec:inference-live}

Après le gel, le calendrier de la saison courante est figé dans \code{staging}. football-data le
complète en mémoire (ADR-0042). Un match du fichier est apparié au calendrier par la clé
\[
  \kappa(m) = (\text{championnat}, \text{saison}, \text{équipe à domicile}, \text{équipe à l'extérieur}),
\]
unique en saison régulière ; la date n'y entre pas, car celle du calendrier figé est souvent fausse.
Une clé ambiguë ou un nom inconnu n'est jamais résolu au hasard. La fraîcheur est propre à chaque
championnat \(L\) :
\[
  C_L = \min\{ d(m) : m \in L,\ m \text{ sans résultat} \} - 1 ,
\]
où \(d(m)\) est la date réelle si elle est connue, la date figée sinon ; une ligne de résultat
inutilisable abaisse \(C_L\) à la veille de sa date. Un match de \(L\) au jour \(J\) a ses variables
d'historique périmées dès que \(C_L < J - 1\). Une répétition sur 2023-24 (gel simulé au 15 février,
\og aujourd'hui \fg{} le 9 mars) rétablit 335 résultats sur 335 sans écart de score, et redonne des
variables de buts et d'Elo identiques à celles du jeu ; seules les variables de tirs changent, l'écart
de log-loss correspondant n'étant pas significatif (ADR-0041).

% ---------------------------------------------------------------------------------------------
\section{Traçabilité}
\label{sec:inference-tracabilite}

Chaque prédiction créée est écrite dans \code{ops.prediction} : match, modèle, horizon, mode
(\code{replay} ou \code{live}, séparés), date de référence, \(\lambda\), loi du total, intervalle et
couverture annoncée, variables utilisées et manquantes, statut, raisons, fraîcheur. La création est
idempotente (une seconde demande renvoie la ligne existante). En live, une prédiction n'est acceptée
qu'avant le coup d'envoi : c'est la condition de l'évaluation prospective de la saison 2026-27.
